% !PRINTED 361
\weberpage{361}

\weq{8}{
\frac{m}{n} = (a_0, a_1, a_2 \ldots a_\nu)
= (a_0, a_1, a_2 \ldots a_\nu, 1).
}

I. \gesperrt{Thus every rational fraction can be developed into a continued fraction which has, at pleasure, an even or an odd number of partial denominators.}

The numbers occurring in the formulae (5),
\[
\frac{m_1}{m_2}, \ \frac{m_2}{m_3}, \ \frac{m_3}{m_4} \ldots,
\]
we call the \gesperrt{closing numbers} of the continued fraction; the numbers $a_0, a_1, a_2 \ldots$ shall be called the \gesperrt{partial denominators} (although $a_0$ does not properly occur as a denominator).

If we replace $m$ by a number $m'$ congruent to it modulo $n$, the two fractions $m:n$ and $m':n$ differ only by an integer; the continued fraction developments therefore differ only in their first terms $a_0$.

\begin{center}
\S.\,116.

\gesperrt{Continued fraction development of irrational numbers.}
\end{center}

If $x$ is a real irrational number, there will always be one and only one integer $a_0$ such that $x$ lies between $a_0$ and $a_0 + 1$. $a_0$ is positive when $x$ is greater than 1, zero when $x$ is a positive proper fraction, and negative when $x$ is negative. If we therefore set
\[
x = a_0 + \frac{1}{x_1},
\]
then $x_1$ is a positive improper fraction; hence, if $a_1$ is the integer lying next \gesperrt{below} $x_1$, $a_1$ is positive. We set
\[
\begin{aligned}
x_1 &= a_1 + \frac{1}{x_2}, \\
x_2 &= a_2 + \frac{1}{x_3}, \\
&\ \cdots\cdots\cdots \\
x_{n-1} &= a_{n-1} + \frac{1}{x_n},
\end{aligned}
\]
%CONT>>
